\section{El movimiento y su descripción}

El movimiento es uno de los fenómenos más familiares de la naturaleza.
Vemos automóviles desplazarse por una carretera, personas caminar, objetos
caer, planetas recorrer sus órbitas y partículas moverse a escalas
microscópicas.

Sin embargo, en Física no basta con afirmar que un cuerpo ``se mueve''.

Para describir un movimiento necesitamos responder preguntas como:

\begin{itemize}
    \item ¿dónde se encuentra el cuerpo?;
    \item ¿respecto de qué estamos indicando su posición?;
    \item ¿cómo cambia esa posición con el tiempo?;
    \item ¿qué camino sigue?;
    \item ¿en qué dirección se desplaza?;
    \item ¿qué distancia recorre?
\end{itemize}

Estas preguntas constituyen el punto de partida de la
\textbf{cinemática}.

\begin{definicion}[Cinemática]
La \textbf{cinemática} es la parte de la mecánica que describe el movimiento
de los cuerpos sin estudiar todavía las causas que producen o modifican ese
movimiento.
\end{definicion}

Por tanto, en cinemática nos interesa describir \textbf{cómo se mueve} un
cuerpo.

Más adelante, en dinámica, estudiaremos \textbf{por qué se mueve de esa
manera}, es decir, analizaremos las fuerzas y las interacciones responsables
del movimiento.

\begin{notapedagogica}
Una forma natural de iniciar la clase es preguntar:

\begin{quote}
Si digo simplemente ``el automóvil se está moviendo'', ¿he descrito
completamente su movimiento?
\end{quote}

La respuesta es no.

Todavía necesitamos saber respecto de qué se mueve, dónde se encuentra,
hacia dónde va y cómo cambia su posición con el tiempo.

La idea que debe quedar establecida desde el comienzo es:

\begin{quote}
Describir un movimiento significa relacionar la posición de un cuerpo con
el tiempo dentro de un sistema de referencia.
\end{quote}
\end{notapedagogica}


\section{Movimiento y reposo son conceptos relativos}

Imaginemos una persona sentada dentro de un autobús que avanza por una
carretera.

Respecto del asiento, la persona permanece en la misma posición.

Podemos decir entonces que está en reposo respecto del autobús.

Sin embargo, respecto de una persona que se encuentra detenida junto a la
carretera, el pasajero cambia continuamente de posición.

Respecto del suelo, el pasajero está en movimiento.

No existe contradicción.

Las dos descripciones utilizan referencias diferentes.

\begin{definicion}[Movimiento]
Un cuerpo se encuentra en \textbf{movimiento} respecto de un sistema de
referencia cuando su posición respecto de ese sistema cambia con el tiempo.
\end{definicion}

\begin{definicion}[Reposo]
Un cuerpo se encuentra en \textbf{reposo} respecto de un sistema de
referencia cuando su posición respecto de ese sistema permanece constante
con el tiempo.
\end{definicion}

\begin{idea}
No tiene sentido físico afirmar únicamente:

\begin{quote}
``este cuerpo está en movimiento''
\end{quote}

sin indicar, de manera explícita o implícita, respecto de qué se está
describiendo el movimiento.
\end{idea}

\begin{ejemplo}[Pasajero dentro de un tren]
Un estudiante permanece sentado dentro de un tren que se desplaza por una
vía recta.

Respecto de su asiento, su posición no cambia. Por tanto, está en reposo.

Respecto de una estación situada junto a la vía, su posición cambia
continuamente. Por tanto, está en movimiento.

El estado de movimiento depende del sistema de referencia elegido.
\end{ejemplo}


\section{Modelo de partícula}

Los cuerpos reales tienen dimensiones.

Un automóvil posee longitud, anchura y altura; una persona tiene brazos,
piernas y puede girar; incluso una pelota puede rotar mientras avanza.

Sin embargo, en muchos problemas no necesitamos describir todos estos
detalles.

Si queremos estudiar el desplazamiento de un automóvil entre dos ciudades,
sus dimensiones son muy pequeñas comparadas con la distancia recorrida.

Podemos entonces representarlo mediante un único punto.

\begin{modelo}[Modelo de partícula]
En el \textbf{modelo de partícula} se representa un cuerpo mediante un punto
al que se atribuyen las propiedades necesarias para estudiar su movimiento,
ignorando su tamaño, forma, orientación y estructura interna cuando estos
aspectos no son relevantes para el problema.
\end{modelo}

Utilizar este modelo no significa afirmar que el objeto realmente sea un
punto.

Significa que, para la pregunta que queremos responder, sus dimensiones no
afectan significativamente la descripción.

\begin{ejemplo}[¿Cuándo podemos considerar un automóvil como partícula?]
Si queremos determinar cuánto tarda un automóvil en viajar de una ciudad a
otra, normalmente podemos representarlo como una partícula.

En cambio, si queremos analizar si el automóvil cabe dentro de un espacio
de estacionamiento, sus dimensiones son esenciales y el modelo de partícula
ya no resulta apropiado.
\end{ejemplo}

\begin{interpretacion}
Un modelo físico conserva los aspectos importantes del fenómeno y omite los
detalles que no son necesarios para el análisis considerado.
\end{interpretacion}


\section{Sistema de referencia}

Para describir dónde se encuentra un cuerpo necesitamos establecer una
referencia.

Decir

\begin{quote}
``el automóvil está a veinte metros''
\end{quote}

no es suficiente.

Debemos preguntar:

\begin{quote}
¿veinte metros desde dónde?
\end{quote}

Para responder necesitamos un \textbf{sistema de referencia}.

\begin{definicion}[Sistema de referencia]
Un \textbf{sistema de referencia} es el conjunto de elementos que utilizamos
para especificar la posición de un cuerpo y describir cómo cambia esa
posición con el tiempo.
\end{definicion}

Para estudiar movimiento en línea recta necesitamos, como mínimo:

\begin{itemize}
    \item un origen;
    \item un eje de coordenadas;
    \item un sentido positivo;
    \item una unidad de longitud;
    \item una medida del tiempo.
\end{itemize}


\subsection{Origen}

El origen es el punto al que asignamos la coordenada

\[ x=0. \]

Su elección es convencional.

Podemos colocar el origen donde resulte más conveniente para resolver el
problema.

\begin{advertencia}
El origen del sistema de coordenadas no tiene que coincidir con la posición
inicial del objeto.

Un cuerpo puede comenzar, por ejemplo, en

\[ x_0=\SI{5}{\meter} \]

o en

\[ x_0=-\SI{3}{\meter}. \]
\end{advertencia}


\subsection{Sentido positivo}

Una vez elegido el eje debemos establecer qué sentido consideraremos
positivo.

Por ejemplo, en una carretera horizontal podemos decidir que:

\[
\text{derecha}\quad\longrightarrow\quad +x
\]

y, por tanto,

\[
\text{izquierda}\quad\longrightarrow\quad -x.
\]

La elección podría hacerse al contrario.

La Física no cambia.

Lo importante es mantener la misma convención durante todo el problema.

\begin{notapedagogica}
Conviene dibujar desde el comienzo un eje en el tablero:

\begin{center}
\begin{tikzpicture}
    \draw[-{Latex},thick] (-5,0) -- (5,0) node[right] {$x$};
    \draw[thick] (0,-0.15) -- (0,0.15);
    \node[below] at (0,-0.15) {$0$};
    \node[above] at (4,0) {sentido positivo};
\end{tikzpicture}
\end{center}

Después puede preguntarse:

\begin{quote}
¿Un número negativo significa necesariamente que el objeto está moviéndose
hacia la izquierda?
\end{quote}

La respuesta es no.

Una coordenada negativa indica \textbf{dónde está} el objeto respecto del
origen, no hacia dónde se está moviendo.
\end{notapedagogica}


\section{Posición}

Una vez definido el sistema de referencia podemos indicar dónde se encuentra
la partícula.

\begin{definicion}[Posición]
La \textbf{posición} de una partícula indica su localización respecto del
sistema de referencia elegido.
\end{definicion}

En movimiento rectilíneo utilizaremos normalmente una coordenada \(x\).

Por ejemplo,

\[ x=\SI{4}{\meter} \]

significa que la partícula se encuentra a cuatro metros del origen en el
sentido positivo del eje.

En cambio,

\[ x=-\SI{4}{\meter} \]

indica que se encuentra a cuatro metros del origen, pero en el sentido
opuesto.

\begin{figure}[H]
    \centering

    \begin{tikzpicture}[scale=1.05]
        \draw[-{Latex},thick] (-5.5,0) -- (5.5,0) node[right] {$x$};

        \foreach \x in {-4,-3,-2,-1,0,1,2,3,4}
        {
            \draw (\x,-0.12) -- (\x,0.12);
        }

        \node[below] at (0,-0.15) {$0$};
        \node[below] at (-4,-0.15) {$-4$};
        \node[below] at (4,-0.15) {$4$};

        \fill (-4,0) circle (2.5pt);
        \node[above] at (-4,0.15) {$A$};

        \fill (4,0) circle (2.5pt);
        \node[above] at (4,0.15) {$B$};

        \node[below] at (0,-0.65) {posición en metros};
    \end{tikzpicture}

    \caption{Dos posiciones distintas respecto del mismo origen.}
\end{figure}

\begin{idea}
El signo de la posición indica en qué lado del origen se encuentra el
objeto.

No indica por sí mismo la dirección en la que se mueve.
\end{idea}


\subsection{Posición como función del tiempo}

Cuando un cuerpo se mueve, su posición puede cambiar conforme transcurre el
tiempo.

Podemos expresar esta dependencia escribiendo

\[ x=x(t). \]

Esta notación se lee:

\begin{quote}
``la posición \(x\) depende del tiempo \(t\)''.
\end{quote}

Por ejemplo, un automóvil podría encontrarse en:

\[
x=\SI{2}{\meter}
\quad\text{cuando}\quad
t=\SI{0}{\second},
\]

y posteriormente en:

\[
x=\SI{10}{\meter}
\quad\text{cuando}\quad
t=\SI{4}{\second}.
\]

La cinemática estudia precisamente cómo cambia la posición con el tiempo.

\begin{interpretacion}
Conocer una única posición nos dice dónde está el objeto en un instante.

Conocer \(x(t)\) nos permite describir cómo evoluciona su posición durante
el movimiento.
\end{interpretacion}


\section{Trayectoria}

Mientras una partícula se mueve ocupa sucesivamente diferentes posiciones.

El conjunto de esas posiciones forma su trayectoria.

\begin{definicion}[Trayectoria]
La \textbf{trayectoria} es el camino geométrico formado por las sucesivas
posiciones que ocupa un cuerpo durante su movimiento.
\end{definicion}

Dependiendo del movimiento, una trayectoria puede ser:

\begin{itemize}
    \item rectilínea;
    \item circular;
    \item parabólica;
    \item curvilínea de otra forma.
\end{itemize}

En este bloque comenzaremos estudiando el caso más sencillo:
el \textbf{movimiento rectilíneo}.

\begin{definicion}[Movimiento rectilíneo]
Un movimiento es \textbf{rectilíneo} cuando la trayectoria del cuerpo es una
línea recta.
\end{definicion}

En este caso basta una sola coordenada, como \(x\), para describir la
posición.


\section{Desplazamiento}

Supongamos que una partícula se encuentra inicialmente en la posición
\(x_i\) y posteriormente en la posición \(x_f\).

Queremos describir el cambio de posición.

\begin{definicion}[Desplazamiento]
El \textbf{desplazamiento} es el cambio de posición de una partícula entre
dos instantes.
\end{definicion}

Para movimiento sobre el eje \(x\),

\[ \boxed{\Delta x=x_f-x_i}. \]

Aquí:

\begin{itemize}
    \item \(x_i\) es la posición inicial;
    \item \(x_f\) es la posición final;
    \item \(\Delta x\) es el desplazamiento.
\end{itemize}

El símbolo \(\Delta\) se lee ``cambio de''.

Por tanto,

\[ \Delta x=x_f-x_i \]

significa:

\begin{quote}
cambio de posición = posición final menos posición inicial.
\end{quote}

La unidad SI del desplazamiento es el metro:

\[ [\Delta x]=\si{\meter}. \]


\subsection{El signo del desplazamiento}

En una dimensión, el signo permite representar el sentido del
desplazamiento.

Si

\[ \Delta x>0, \]

el desplazamiento ocurre en el sentido positivo del eje.

Si

\[ \Delta x<0, \]

el desplazamiento ocurre en el sentido negativo.

Si

\[ \Delta x=0, \]

la posición final coincide con la posición inicial.

\begin{advertencia}
Un desplazamiento igual a cero no significa necesariamente que el objeto no
se movió.

Puede haberse alejado de su posición inicial y después haber regresado a
ella.
\end{advertencia}


\begin{ejemplo}[Desplazamiento positivo]
Un estudiante camina en línea recta desde

\[ x_i=-\SI{2}{\meter} \]

hasta

\[ x_f=\SI{5}{\meter}. \]

El desplazamiento es

\[
\Delta x=x_f-x_i
=\SI{5}{\meter}-(-\SI{2}{\meter})
=\SI{7}{\meter}.
\]

Por tanto,

\[ \boxed{\Delta x=+\SI{7}{\meter}}. \]

El signo positivo indica que el cambio de posición ocurrió en el sentido
positivo del eje.
\end{ejemplo}


\begin{ejemplo}[Desplazamiento negativo]
Una persona se encuentra inicialmente en

\[ x_i=\SI{6}{\meter} \]

y termina en

\[ x_f=-\SI{3}{\meter}. \]

Entonces,

\[
\Delta x=x_f-x_i
=-\SI{3}{\meter}-\SI{6}{\meter}
=-\SI{9}{\meter}.
\]

Por tanto,

\[ \boxed{\Delta x=-\SI{9}{\meter}}. \]

La magnitud del desplazamiento es

\[ \abs{\Delta x}=\SI{9}{\meter}, \]

y el signo negativo indica que el desplazamiento neto fue en el sentido
negativo del eje.
\end{ejemplo}


\section{Distancia recorrida}

El desplazamiento solamente compara la posición inicial con la posición
final.

Sin embargo, en ocasiones nos interesa conocer la longitud total del camino
recorrido.

Esta magnitud es la distancia.

\begin{definicion}[Distancia recorrida]
La \textbf{distancia recorrida} es la longitud total de la trayectoria
seguida por un cuerpo durante un intervalo de movimiento.
\end{definicion}

Representaremos la distancia recorrida mediante \(s\).

La distancia es una magnitud escalar y cumple

\[ s\geq 0. \]

Su unidad SI es también el metro.

\begin{idea}
La distancia responde:

\begin{quote}
¿Cuánto camino recorrió?
\end{quote}

El desplazamiento responde:

\begin{quote}
¿Cuánto cambió su posición y en qué sentido?
\end{quote}
\end{idea}


\section{Distancia y desplazamiento no son lo mismo}

Esta diferencia es una de las ideas más importantes de la cinemática.

Supongamos que una persona sale del origen, camina hasta

\[ x=\SI{6}{\meter}, \]

y después regresa hasta

\[ x=-\SI{2}{\meter}. \]

La trayectoria puede representarse como:

\begin{center}
\begin{tikzpicture}[scale=0.8]
    \draw[-{Latex},thick] (-4,0) -- (8,0) node[right] {$x$};

    \foreach \x in {-2,0,2,4,6}
    {
        \draw (\x,-0.12) -- (\x,0.12);
        \node[below] at (\x,-0.15) {\(\x\)};
    }

    \draw[-{Latex},thick] (0,0.55) -- (6,0.55);
    \draw[-{Latex},thick] (6,0.9) -- (-2,0.9);

    \fill (0,0) circle (2.5pt);
    \fill (6,0) circle (2.5pt);
    \fill (-2,0) circle (2.5pt);

    \node[above] at (0,1.15) {inicio};
    \node[above] at (-2,1.15) {final};
\end{tikzpicture}
\end{center}

En el primer tramo recorre

\[ \SI{6}{\meter}. \]

En el segundo tramo, desde \(6\) hasta \(-2\), recorre

\[ \SI{8}{\meter}. \]

Por tanto, la distancia total es

\[
s=\SI{6}{\meter}+\SI{8}{\meter}
=\SI{14}{\meter}.
\]

En cambio, el desplazamiento es

\[
\Delta x=x_f-x_i
=-\SI{2}{\meter}-\SI{0}{\meter}
=-\SI{2}{\meter}.
\]

Por tanto,

\[
\boxed{s=\SI{14}{\meter}},
\qquad
\boxed{\Delta x=-\SI{2}{\meter}}.
\]

El cuerpo recorrió catorce metros, pero terminó solamente dos metros a la
izquierda de su posición inicial.


\subsection{Relación entre distancia y desplazamiento}

Siempre se cumple

\[ \boxed{s\geq\abs{\Delta x}}. \]

La igualdad

\[ s=\abs{\Delta x} \]

se cumple cuando el cuerpo se desplaza entre las posiciones inicial y final
sin invertir el sentido de su movimiento a lo largo de la trayectoria
rectilínea.

Si cambia de sentido durante el recorrido, generalmente

\[ s>\abs{\Delta x}. \]

\begin{ejemplo}[Regreso al punto de partida]
Un corredor avanza \(\SI{100}{\meter}\) por una pista recta y luego regresa
\(\SI{100}{\meter}\) hasta el punto de partida.

La distancia recorrida es

\[
s=\SI{100}{\meter}+\SI{100}{\meter}
=\SI{200}{\meter}.
\]

Sin embargo,

\[ x_f=x_i, \]

por lo que

\[
\Delta x=x_f-x_i=0.
\]

Así,

\[ \boxed{s=\SI{200}{\meter}}, \qquad \boxed{\Delta x=0}. \]

Este ejemplo muestra con claridad que un objeto puede recorrer una distancia
considerable y tener desplazamiento nulo.
\end{ejemplo}


\section{Instante e intervalo de tiempo}

Para describir el movimiento no basta conocer dónde se encuentra el objeto.

También necesitamos saber \textbf{cuándo} ocupa cada posición.

Representaremos el tiempo mediante \(t\).

Un valor como

\[ t=\SI{4}{\second} \]

representa un instante determinado respecto del origen temporal elegido.

Si un fenómeno comienza en el instante \(t_i\) y termina en \(t_f\), el
intervalo de tiempo transcurrido es

\[ \boxed{\Delta t=t_f-t_i}. \]

Para un proceso físico en el que \(t_f>t_i\),

\[ \Delta t>0. \]

La unidad SI del tiempo es el segundo:

\[ [t]=[\Delta t]=\si{\second}. \]

\begin{advertencia}
No deben confundirse el instante \(t\) y el intervalo de tiempo
\(\Delta t\).

Por ejemplo, un evento puede comenzar a

\[ t_i=\SI{3}{\second} \]

y terminar a

\[ t_f=\SI{8}{\second}. \]

El intervalo transcurrido no es \(\SI{8}{\second}\), sino

\[
\Delta t=\SI{8}{\second}-\SI{3}{\second}
=\SI{5}{\second}.
\]
\end{advertencia}


\section{Rapidez y velocidad: primera aproximación}

Hasta ahora podemos indicar dónde está un cuerpo y cuánto cambia su
posición.

Sin embargo, aún falta una pregunta fundamental:

\begin{quote}
¿Qué tan rápido ocurre ese cambio?
\end{quote}

Esto conduce a dos conceptos relacionados pero diferentes:
\textbf{rapidez} y \textbf{velocidad}.


\subsection{Rapidez}

La rapidez indica qué tan rápidamente se recorre distancia.

Es una magnitud escalar.

Por tanto, una rapidez como

\[ \SI{20}{\meter\per\second} \]

no contiene por sí sola información sobre la dirección del movimiento.

\begin{interpretacion}
Dos automóviles pueden tener la misma rapidez y estar desplazándose en
sentidos opuestos.
\end{interpretacion}


\subsection{Velocidad}

La velocidad describe qué tan rápidamente cambia la posición y además
incorpora información sobre la dirección del movimiento.

La velocidad es una magnitud vectorial.

En movimiento rectilíneo podemos representar su sentido mediante un signo.

Si hemos elegido la derecha como sentido positivo, entonces una velocidad
positiva corresponde al movimiento hacia la derecha y una velocidad
negativa al movimiento hacia la izquierda.

\begin{idea}
\textbf{Rapidez} y \textbf{velocidad} no son sinónimos en Física.

La rapidez es escalar.

La velocidad es vectorial.
\end{idea}

\begin{ejemplo}[Misma rapidez, distinta velocidad]
Dos automóviles se desplazan por una carretera recta.

El automóvil A se mueve hacia la derecha con rapidez

\[ \SI{20}{\meter\per\second}. \]

El automóvil B se mueve hacia la izquierda con la misma rapidez.

Si elegimos la derecha como sentido positivo, podemos representar sus
velocidades como

\[ v_A=+\SI{20}{\meter\per\second}, \]

\[ v_B=-\SI{20}{\meter\per\second}. \]

Ambos poseen la misma rapidez:

\[ \SI{20}{\meter\per\second}, \]

pero sus velocidades son diferentes porque se mueven en sentidos opuestos.
\end{ejemplo}

\begin{notapedagogica}
En este punto conviene introducir solamente la diferencia conceptual entre
rapidez y velocidad.

Las definiciones matemáticas de rapidez media, velocidad media y velocidad
instantánea se desarrollarán en el capítulo siguiente.

Así se evita convertir esta primera clase en una colección de fórmulas antes
de haber establecido correctamente el lenguaje del movimiento.
\end{notapedagogica}


\section{Posición, desplazamiento y velocidad son conceptos diferentes}

En movimiento rectilíneo aparecen varios signos que pueden resultar
confusos si no distinguimos qué representa cada magnitud.

Supongamos que hemos elegido la derecha como sentido positivo.

Un cuerpo puede encontrarse en

\[ x=-\SI{5}{\meter} \]

y moverse hacia la derecha.

En ese instante:

\[
x<0,
\]

pero su velocidad puede ser positiva:

\[
v>0.
\]

También podría encontrarse en

\[ x=+\SI{5}{\meter} \]

y moverse hacia la izquierda.

Entonces:

\[
x>0,
\]

pero

\[
v<0.
\]

\begin{idea}
El signo de la posición responde:

\begin{quote}
¿en qué lado del origen está el cuerpo?
\end{quote}

El signo del desplazamiento responde:

\begin{quote}
¿en qué sentido ocurrió el cambio neto de posición?
\end{quote}

El signo de la velocidad responde:

\begin{quote}
¿en qué sentido se está moviendo?
\end{quote}

Son preguntas diferentes.
\end{idea}


\section{¿Qué información necesitamos para describir un movimiento rectilíneo?}

Después de introducir los conceptos anteriores podemos resumir la estructura
de una descripción cinemática.

Para describir un movimiento en línea recta debemos especificar:

\begin{enumerate}
    \item el sistema de referencia;
    \item el origen;
    \item el eje sobre el que ocurre el movimiento;
    \item el sentido positivo;
    \item la posición del objeto;
    \item cómo cambia esa posición con el tiempo.
\end{enumerate}

A partir de esta información podremos construir posteriormente cantidades
como velocidad y aceleración.

\begin{notapedagogica}
Una frase útil para cerrar esta parte de la explicación es:

\begin{quote}
En cinemática no comenzamos preguntando qué fuerza actúa sobre el objeto.

Primero construimos una descripción precisa de dónde está y cómo cambia su
posición con el tiempo.
\end{quote}
\end{notapedagogica}


\section{Ejemplo integrador}

\begin{ejemplo}[Movimiento sobre una carretera recta]
Una persona se mueve sobre una carretera recta.

Elegimos un árbol como origen y definimos hacia el este como sentido
positivo.

La persona comienza en

\[ x_i=-\SI{4}{\meter}. \]

Camina primero hasta

\[ x=\SI{7}{\meter}, \]

y posteriormente regresa hasta

\[ x_f=\SI{2}{\meter}. \]

Determine:

\begin{enumerate}[label=\alph*)]
    \item qué significa la posición inicial;
    \item la distancia total recorrida;
    \item el desplazamiento;
    \item el sentido del desplazamiento neto.
\end{enumerate}
\end{ejemplo}

La posición inicial

\[ x_i=-\SI{4}{\meter} \]

significa que la persona comienza a cuatro metros al oeste del árbol.

En el primer tramo pasa de \(-4\) a \(7\):

\[
s_1=\abs{7-(-4)}\,\si{\meter}
=\SI{11}{\meter}.
\]

En el segundo tramo pasa de \(7\) a \(2\):

\[
s_2=\abs{2-7}\,\si{\meter}
=\SI{5}{\meter}.
\]

Por tanto, la distancia total es

\[
s=s_1+s_2
=\SI{11}{\meter}+\SI{5}{\meter}
=\SI{16}{\meter}.
\]

El desplazamiento depende únicamente de las posiciones inicial y final:

\[
\Delta x=x_f-x_i
=\SI{2}{\meter}-(-\SI{4}{\meter})
=\SI{6}{\meter}.
\]

Por tanto,

\[ \boxed{s=\SI{16}{\meter}} \]

y

\[ \boxed{\Delta x=+\SI{6}{\meter}}. \]

El signo positivo del desplazamiento indica que la posición final está seis
metros en el sentido positivo respecto de la posición inicial.

\begin{interpretacion}
La persona recorrió dieciséis metros, pero su cambio neto de posición fue
solamente seis metros hacia el este.

Esto muestra nuevamente que distancia y desplazamiento describen aspectos
diferentes del mismo movimiento.
\end{interpretacion}


\section{Errores y confusiones frecuentes}

\begin{errorfrecuente}
\textbf{Pensar que movimiento y reposo son absolutos.}

Siempre deben establecerse respecto de un sistema de referencia.
\end{errorfrecuente}

\begin{errorfrecuente}
\textbf{Confundir posición con distancia recorrida.}

La posición indica dónde está el objeto respecto del origen.

La distancia indica cuánto camino ha recorrido.
\end{errorfrecuente}

\begin{errorfrecuente}
\textbf{Pensar que una posición negativa significa que el objeto se mueve
hacia atrás.}

Una posición negativa únicamente indica que el objeto está situado en el
lado negativo del origen.
\end{errorfrecuente}

\begin{errorfrecuente}
\textbf{Calcular siempre la distancia mediante}

\[ \abs{x_f-x_i}. \]

Esta expresión proporciona la magnitud del desplazamiento.

Solo coincide con la distancia recorrida cuando el cuerpo no cambia de
sentido durante el trayecto rectilíneo considerado.
\end{errorfrecuente}

\begin{errorfrecuente}
\textbf{Pensar que un desplazamiento cero significa que no hubo movimiento.}

Un cuerpo puede recorrer una trayectoria y regresar a su posición inicial.

En ese caso la distancia es distinta de cero, pero el desplazamiento es cero.
\end{errorfrecuente}

\begin{errorfrecuente}
\textbf{Confundir rapidez y velocidad.}

La rapidez es escalar.

La velocidad contiene información sobre la dirección y el sentido del
movimiento.
\end{errorfrecuente}

\begin{errorfrecuente}
\textbf{Confundir el origen con la posición inicial.}

El origen se elige como parte del sistema de referencia.

La partícula puede iniciar su movimiento en cualquier coordenada.
\end{errorfrecuente}


\section{Secuencia conceptual para explicar el tema}

\begin{notapedagogica}
Una secuencia clara para dictar esta clase es:

\begin{enumerate}

    \item Comenzar preguntando qué significa afirmar que un cuerpo se mueve.

    \item Utilizar el ejemplo de un pasajero dentro de un autobús o tren para
    mostrar que movimiento y reposo dependen del sistema de referencia.

    \item Introducir la cinemática como descripción del movimiento sin
    estudiar todavía sus causas.

    \item Explicar el modelo de partícula y por qué simplificamos los cuerpos.

    \item Dibujar un eje horizontal en el tablero.

    \item Elegir un origen y un sentido positivo.

    \item Colocar varios objetos en posiciones positivas y negativas.

    \item Preguntar explícitamente si una coordenada negativa indica una
    dirección de movimiento.

    \item Definir posición.

    \item Mostrar que durante el movimiento la posición depende del tiempo:

    \[
    x=x(t).
    \]

    \item Introducir el concepto de trayectoria.

    \item Particularizar al movimiento rectilíneo.

    \item Definir desplazamiento:

    \[
    \Delta x=x_f-x_i.
    \]

    \item Interpretar cuidadosamente su signo.

    \item Definir distancia recorrida.

    \item Resolver inmediatamente un ejemplo con cambio de sentido para
    diferenciar distancia y desplazamiento.

    \item Introducir instante e intervalo de tiempo.

    \item Presentar cualitativamente rapidez y velocidad.

    \item Utilizar un ejemplo de dos objetos con la misma rapidez y sentidos
    opuestos.

    \item Terminar estableciendo qué información se necesita para describir
    completamente un movimiento rectilíneo.

\end{enumerate}
\end{notapedagogica}


\section{Preguntas para comprobar comprensión}

Antes de avanzar al estudio cuantitativo de la velocidad pueden plantearse
preguntas como:

\begin{enumerate}

    \item ¿Puede un cuerpo estar en reposo respecto de un observador y en
    movimiento respecto de otro?

    \item ¿Qué elementos debemos definir antes de asignar una posición a un
    objeto?

    \item ¿El origen debe coincidir necesariamente con la posición inicial?

    \item Si un objeto está en \(x=-\SI{5}{\meter}\), ¿podemos saber hacia
    dónde se mueve?

    \item ¿Qué diferencia existe entre posición y desplazamiento?

    \item ¿Qué significa físicamente un desplazamiento negativo?

    \item ¿Puede la distancia recorrida ser negativa?

    \item ¿Puede el desplazamiento ser negativo?

    \item ¿Puede un cuerpo recorrer una distancia distinta de cero y tener
    desplazamiento cero?

    \item ¿En qué situación se cumple

    \[
    s=\abs{\Delta x}?
    \]

    \item Si una persona camina diez metros hacia la derecha y después diez
    metros hacia la izquierda, ¿cuál es la distancia recorrida y cuál es su
    desplazamiento?

    \item ¿Cuál es la diferencia conceptual entre rapidez y velocidad?

    \item ¿Dos objetos pueden tener la misma rapidez y velocidades
    diferentes?

\end{enumerate}


\section{Conexión con el tema siguiente}

Hasta ahora hemos aprendido a establecer:

\begin{itemize}
    \item dónde se encuentra un objeto;
    \item cuánto cambia su posición;
    \item cuánto camino recorre;
    \item cómo representar el sentido de un movimiento rectilíneo.
\end{itemize}

La siguiente pregunta será:

\begin{quote}
¿Con qué rapidez cambia la posición?
\end{quote}

Para responderla introduciremos formalmente la
\textbf{velocidad media}, la \textbf{rapidez media} y posteriormente la
\textbf{velocidad instantánea}.

También aprenderemos a representar el movimiento mediante gráficas de
posición en función del tiempo.


\begin{resumen}

\begin{itemize}

    \item La cinemática describe el movimiento sin estudiar inicialmente sus
    causas.

    \item Movimiento y reposo deben definirse respecto de un sistema de
    referencia.

    \item El modelo de partícula permite representar un cuerpo como un punto
    cuando sus dimensiones no son relevantes para el problema.

    \item Para describir movimiento en línea recta elegimos un origen, un eje
    y un sentido positivo.

    \item La posición \(x\) indica dónde se encuentra una partícula respecto
    del origen.

    \item Una posición negativa no indica necesariamente que el cuerpo se
    mueva en el sentido negativo.

    \item La posición de un cuerpo en movimiento puede expresarse como

    \[
    x=x(t).
    \]

    \item La trayectoria es el camino formado por las sucesivas posiciones
    del cuerpo.

    \item En movimiento rectilíneo la trayectoria es una línea recta.

    \item El desplazamiento es el cambio de posición:

    \[
    \boxed{\Delta x=x_f-x_i}.
    \]

    \item El signo del desplazamiento indica su sentido respecto del eje
    elegido.

    \item La distancia recorrida \(s\) es la longitud total de la trayectoria
    y nunca es negativa.

    \item En general,

    \[
    \boxed{s\geq\abs{\Delta x}}.
    \]

    \item Distancia y desplazamiento solo tienen igual magnitud en un
    movimiento rectilíneo sin cambio de sentido entre las posiciones
    consideradas.

    \item Un objeto puede recorrer una distancia distinta de cero y terminar
    con desplazamiento cero.

    \item El intervalo de tiempo es

    \[
    \boxed{\Delta t=t_f-t_i}.
    \]

    \item La rapidez es una magnitud escalar.

    \item La velocidad es una magnitud vectorial y, en movimiento
    unidimensional, su sentido puede representarse mediante un signo.

\end{itemize}

\end{resumen}